\documentclass[conference]{IEEEtran}
\IEEEoverridecommandlockouts
% ---------------------------------------------------------------------------
% Versao com linguagem e estrutura segundo as normas da ABNT (NBR 6022:
% artigo; NBR 6028: resumo; NBR 10520: citacoes; NBR 6023: referencias;
% NBR 6024: numeracao progressiva), aplicadas sobre o "IEEE Conference
% Template" oficial do Overleaf, que a especificacao da disciplina exige
% (leiaute em duas colunas, fontes e margens do template).
% Para compilar: use o template no Overleaf, substitua o main.tex por este
% arquivo e envie docs/figuras/capturas_sistema.png e
% docs/figuras/custo_vs_carga.png para a pasta do projeto.
% ---------------------------------------------------------------------------
\usepackage[utf8]{inputenc}
\usepackage[T1]{fontenc}
\usepackage[english,brazil]{babel}
\usepackage{amsmath}
\usepackage{amssymb}
\usepackage{graphicx}
\usepackage{textcomp}
\usepackage{xcolor}
\usepackage{booktabs}
\usepackage{url}

% Numeracao progressiva em algarismos arabicos (NBR 6024)
\renewcommand{\thesection}{\arabic{section}}
\renewcommand{\thesubsection}{\thesection.\arabic{subsection}}
\renewcommand{\thesubsubsection}{\thesubsection.\arabic{subsubsection}}
\renewcommand{\thesectiondis}{\thesection}
\renewcommand{\thesubsectiondis}{\thesection.\arabic{subsection}}
\renewcommand{\thesubsubsectiondis}{\thesubsection.\arabic{subsubsection}}
\renewcommand{\thetable}{\arabic{table}}
\renewcommand{\thefigure}{\arabic{figure}}

\newcommand{\fonte}[1]{\par\vspace{2pt}{\footnotesize Fonte: #1\par}}

% Legenda ABNT (NBR 14724): "Tabela 1 -- Titulo", em uma linha, sobre o objeto
\makeatletter
\long\def\@makecaption#1#2{%
  \par\addvspace{4pt}%
  \parbox[t]{\hsize}{\centering\normalfont\footnotesize #1 \textendash\ #2}%
  \par\addvspace{4pt}}
\makeatother

\begin{document}

\title{Visualização Interativa de Elevação Costeira com Nível de Detalhe
Adaptativo e Simulação Oceânica para Análise de Cenários de Subida do
Nível do Mar}

\author{
\IEEEauthorblockN{Guilherme Augusto Scaglia, João Pedro Denardo, Pedro Henrique Oliveira de Souza}
\IEEEauthorblockA{
Engenharia de Computação\\
Fundação Hermínio Ometto (FHO)\\
Araras, SP, Brasil\\
\{scaglia, denardo749, pedro1204\}@alunos.fho.edu.br}
}

\maketitle

\begin{abstract}
Este trabalho apresenta um sistema interativo de visualização científica
para a análise de cenários de erosão costeira e de subida do nível do mar,
implementado sobre um Modelo Digital de Elevação (DEM) real da Baixada
Santista (SP). Os dados combinam observações de radar orbital e batimetria,
isto é, a profundidade do fundo do mar. O objetivo foi determinar, por meio
de experimento, a partir de qual carga geométrica a técnica de nível de
detalhe adaptativo CDLOD (\emph{Continuous Distance-Dependent Level of
Detail}) torna-se indispensável para manter a interatividade da aplicação em uma GPU de computador de mesa. A superfície oceânica
foi modelada por seis ondas de Gerstner, com altura de base ajustável em
tempo real, o que permite visualizar cenários de subida do mar de 0, +0,5,
+1 e +2 metros. Um teste de desempenho automatizado, sem janela gráfica,
executado em uma unidade de processamento gráfico (GPU) NVIDIA RTX 3060,
comparou a renderização com e sem nível de detalhe em cinco resoluções de
terreno. Os resultados indicam um ponto de virada entre 2,7 e 4,2 milhões de
triângulos, conforme o critério (percentil 95 ou média), a partir do qual
o CDLOD passa a ser mais rápido do que a malha completa. Conclui-se que o nível de detalhe adaptativo é o que
torna viável a visualização tridimensional de inundação costeira em tempo
real.

Palavras-chave: Visualização científica. Nível de detalhe. Renderização de
terreno. Simulação de oceano. Subida do nível do mar.

\textit{Interactive Coastal Elevation Visualization with Adaptive Level of
Detail and Ocean Simulation for Sea Level Rise Scenario Analysis}

\textit{Abstract}---This paper presents an interactive scientific
visualization system for analyzing coastal erosion and sea level rise
scenarios, built on a real Digital Elevation Model (DEM) of the Baixada
Santista region (Brazil). The data combine orbital radar observations and
bathymetry, that is, the depth of the seafloor. The goal was to determine
experimentally the geometric load beyond which the adaptive level of detail
technique CDLOD (Continuous Distance-Dependent Level of Detail) becomes
indispensable to keep the application interactive on the tested desktop GPU. The ocean surface was modeled with
six Gerstner waves, with a base height adjustable in real time, which
allows sea level rise scenarios of 0, +0.5, +1 and +2 meters to be
visualized. An automated windowless benchmark, run on an NVIDIA RTX 3060
graphics processing unit (GPU), compared rendering with and without level
of detail at five terrain resolutions. The results indicate a tipping point
between 2.7 and 4.2 million triangles, depending on the criterion (95th
percentile or mean), beyond which CDLOD becomes faster than the full mesh. It is concluded that adaptive level of detail is what
makes real-time three-dimensional visualization of coastal flooding
feasible.

Keywords: Scientific visualization. Level of detail. Terrain rendering.
Ocean simulation. Sea level rise.
\end{abstract}

\section{Introdução}
\label{sec:intro}

O aumento do nível médio dos oceanos, associado às mudanças climáticas de
origem antrópica, figura entre os fenômenos ambientais de maior impacto
sobre as regiões costeiras. Esse processo intensifica a erosão das linhas
de costa, amplia a vulnerabilidade das áreas urbanas litorâneas a
inundações e representa risco direto a infraestruturas, ecossistemas e
populações. A previsão e a comunicação visual desses riscos são, portanto,
centrais tanto para o planejamento urbano e ambiental quanto para a
conscientização pública.

Segundo Brown et al. (2006), um dos principais entraves à comunicação do
risco costeiro à população e aos gestores públicos é a baixa capacidade
das ferramentas tradicionais de geovisualização de transmitir a dinâmica
do fenômeno. A maioria delas adota o chamado modelo de banheira
(\emph{bathtub model}), no qual toda área abaixo de determinada cota de
elevação é simplesmente marcada como inundada, sem simulação da
superfície da água. Para mitigar essa limitação, este trabalho integra a
renderização de terreno em tempo real e a simulação procedural de ondas
oceânicas sobre dados reais de relevo (topografia) e de profundidade do
fundo do mar (batimetria).

A contribuição técnica deste trabalho precisa ser delimitada. O nível de
detalhe adaptativo (LOD, do inglês \emph{level of detail}) é tratado como
meio, e não como fim: não se propõe uma nova técnica de LOD. A pergunta de
pesquisa que orienta a avaliação experimental (Seção~\ref{sec:resultados})
não é se a técnica CDLOD (\emph{Continuous Distance-Dependent Level of
Detail}) funciona, o que já foi estabelecido por Strugar (2009),
mas a partir de que carga geométrica, em um cenário real de visualização
de inundação costeira, o CDLOD deixa de ser um refinamento de desempenho e
passa a ser condição necessária para a interatividade. A questão é relevante porque o ponto em que isso ocorre depende da unidade de processamento gráfico (GPU) utilizada, e a medição foi feita em uma única GPU.

O objetivo geral é construir uma ferramenta interativa de visualização de
cenários de subida do nível do mar sobre um Modelo Digital de Elevação
(DEM) real e determinar experimentalmente a partir de qual carga geométrica
o LOD adaptativo se torna necessário. Os objetivos específicos são: a)
implementar o terreno com CDLOD sobre o DEM; b) simular a superfície
oceânica com altura de base ajustável; e c) medir o tempo de quadro e o
consumo de memória com e sem LOD. A hipótese é que, acima de certa carga
geométrica, a malha completa deixa de ser interativa, enquanto o CDLOD
mantém custo aproximadamente constante.

Diferentemente dos sistemas de geovisualização baseados em Sistemas de
Informação Geográfica (SIG), que representam a subida do mar por camadas
estáticas sobrepostas a um mapa bidimensional, propõe-se uma cena
tridimensional interativa, construída sobre o DEM. Nela, o
observador navega pelo relevo costeiro e ajusta, em tempo real, a altura
da superfície oceânica simulada, observando quais áreas seriam submersas
em cada cenário. A ferramenta destina-se ao público leigo e a gestores
ambientais.

O sistema foi desenvolvido em Python, com a biblioteca ModernGL para
acesso à interface de programação (API) gráfica OpenGL, a GLFW para o gerenciamento de janela e de
entradas e a PyGLM para a álgebra linear das transformações de câmera e
de projeção. A ferramenta busca não apenas renderizar um terreno
convincente, mas também servir como instrumento de análise visual de
baixo custo computacional, complementar às abordagens baseadas em SIG.

Este artigo está organizado da seguinte forma. A
Seção~\ref{sec:fund} apresenta a fundamentação teórica e os trabalhos
relacionados. A Seção~\ref{sec:metod} descreve a metodologia e a
modelagem computacional. A Seção~\ref{sec:resultados} apresenta os
resultados da avaliação de desempenho. A Seção~\ref{sec:conclusao}
traz a conclusão, as limitações e os trabalhos futuros.

\subsection{Justificativa da Escolha do Cenário e das Técnicas}

A escolha do cenário de erosão costeira justifica-se por três razões.
Primeiro, trata-se de um cenário com embasamento em trabalhos de
geovisualização científica (BROWN et al., 2006; DIPAOLA; BHARDWAJ; SAM,
2023). Segundo, o parâmetro central da simulação de água, a altura de base
ajustável, corresponde a uma variável de interesse real: as projeções de
subida do mar, expressas em metros. Terceiro, o problema impõe uma escala
compatível com o escopo do projeto. Ao contrário da renderização de um
planeta inteiro (CLASEN; HEGE, 2006), um recorte costeiro local, a Baixada
Santista, detalhada na Subseção~\ref{sub:aquisicao}, permite simplificar a
topologia do terreno de esférica para plana sem descaracterizar o
problema.

A escolha das técnicas, CDLOD para o terreno e ondas de Gerstner para a
água, segue o mesmo raciocínio de custo e benefício, confirmado
experimentalmente na Seção~\ref{sec:resultados}. O CDLOD dispensa
estruturas de dados hierárquicas mais complexas do que as de outras
técnicas de LOD. O modelo de Gerstner dispensa o custo das transformadas de
Fourier exigidas por modelos espectrais de oceano, mantendo fidelidade
visual suficiente para a comunicação de risco e sem inviabilizar a
execução em hardware de menor capacidade, que é o objetivo final da
aplicação.

\section{Fundamentação Teórica e Trabalhos Relacionados}
\label{sec:fund}

A revisão bibliográfica abrange sete trabalhos, organizados em três
frentes: a renderização eficiente de terrenos extensos, a simulação de
ondas oceânicas e a geovisualização de risco costeiro. As
Subseções~\ref{sub:terreno} a~\ref{sub:geovis} descrevem o que cada frente
propõe e quais são suas limitações para o problema tratado. A
Subseção~\ref{sub:sintese} reúne essas observações e explicita o que este
trabalho adota de cada uma e no que se diferencia.

\subsection{Renderização de Terreno em Larga Escala}
\label{sub:terreno}

Um terreno real extenso, representado como malha de triângulos, pode
conter milhões de triângulos, o que torna inviável desenhá-lo por inteiro a
cada quadro. A solução usual é o nível de detalhe (LOD): usar muitos
triângulos nas regiões próximas ao observador, onde o detalhe é
perceptível, e poucos nas regiões distantes.

Tanner, Migdal e Jones (1998) introduziram o \emph{clipmap}, estrutura que
mantém em memória uma ``janela'' de dados em alta resolução centrada no
observador, cercada por janelas sucessivamente maiores e de resolução
menor. Assim, apenas os dados próximos precisam ser armazenados em alta
resolução, o que permite renderizar mapas de elevação maiores do que a
memória da GPU comportaria. Clasen e Hege (2006) estenderam essa ideia a
terrenos esféricos, com anéis circulares de detalhe, de modo a visualizar
planetas inteiros a partir de modelos digitais de elevação de alta
resolução.

Ambos os trabalhos visam grandes extensões, até a escala planetária. Para
um recorte costeiro local, de superfície aproximadamente plana, a
parametrização esférica e a manutenção de janelas aninhadas de dados
acrescentam uma complexidade que o problema não exige.

\subsection{Nível de Detalhe Contínuo Baseado em Distância}
\label{sub:cdlod-rel}

Strugar (2009) propôs o CDLOD (\emph{Continuous Distance-Dependent Level of
Detail}). O terreno é organizado em uma \emph{quadtree}, estrutura em
árvore na qual um bloco quadrado pode ser dividido em quatro blocos
menores, e assim recursivamente. Blocos próximos ao observador são
subdivididos e ganham detalhe, enquanto blocos distantes permanecem grandes
e grosseiros. A decisão depende de uma função contínua da distância, e a
troca de nível é suavizada por \emph{geomorphing}: os vértices deslizam
gradualmente de uma resolução para a seguinte, o que evita saltos visuais
abruptos e fendas entre blocos vizinhos.

Em comparação com o \emph{clipmap}, o CDLOD é mais simples de implementar,
pois dispensa janelas de dados aninhadas e se resume a escolher, a cada
quadro, quais blocos desenhar. O trabalho original foi proposto e avaliado
para terrenos genéricos, sem vínculo com um cenário de aplicação
específico.

\subsection{Simulação Procedural de Ondas Oceânicas}
\label{sub:ondas-rel}

Fournier e Reeves (1986) propuseram um dos primeiros modelos procedurais de
ondas para síntese de imagens, baseado no modelo de Gerstner. Nele, cada
ponto da superfície da água descreve uma trajetória circular ou elíptica, o
que produz cristas estreitas e vales largos, como em ondas reais. Somando
várias ondas de tamanhos diferentes, obtém-se uma superfície de aspecto
natural, sem exigir as transformadas de Fourier dos modelos espectrais. O trabalho
considera ainda a influência da batimetria (profundidade do fundo do mar)
sobre a refração e a arrebentação das ondas próximas à costa. Welch (2019)
implementa uma simulação de oceano em tempo real com ondas de Gerstner
calculadas em \emph{shaders} GLSL sobre uma geometria estática em
\emph{clipmap}: a GPU desloca os vértices e calcula as normais a cada quadro.
O trabalho acrescenta iluminação fisicamente baseada, reflexos e espuma e
relata taxa média superior a 60 quadros por segundo em hardware de gama
intermediária, a 1920$\times$1080. Transformada rápida de Fourier, refração
e águas rasas ficaram como trabalho futuro.

Nenhum desses trabalhos trata o nível médio da água como variável de
cenário: o oceano é apresentado em nível fixo.

\subsection{Geovisualização de Erosão Costeira e Subida do Nível do Mar}
\label{sub:geovis}

Brown et al. (2006) discutem a importância de visualizar linhas de costa em
evolução como ferramenta de conscientização e planejamento frente à erosão,
e propõem uma metodologia que integra a modelagem científica da erosão a
modelos digitais de elevação (DEM) e a sistemas de informação geográfica.
De forma semelhante, DiPaola, Bhardwaj e Sam (2023) constroem mapas
interativos de subida do nível do mar sobre dados topográficos reais
(levantamento aéreo de alta resolução por LiDAR, do inglês \emph{Light
Detection and Ranging}) do North Shore de Vancouver, Canadá,
quantificando a porcentagem de área afetada em diferentes cenários.

Esses trabalhos são cientificamente rigorosos, mas apresentam o fenômeno por
meio de mapas bidimensionais ou de modelos de inundação estáticos, sem
simular a superfície da água.

\subsection{Síntese e Posicionamento deste Trabalho}
\label{sub:sintese}

Da primeira frente, mantém-se a ideia de concentrar o detalhe próximo ao
observador. Como o recorte de interesse é local e plano, adota-se a
\emph{quadtree} do CDLOD em lugar da parametrização esférica do
\emph{clipmap}. Da segunda frente, adota-se o modelo de Gerstner, na
formulação de Fournier e Reeves (1986), cuja viabilidade em tempo real na GPU é demonstrada por Welch
(2019), com uma diferença: a altura média
da água é um parâmetro ajustável em tempo real, usado para representar
diferentes cenários de subida do mar. A terceira frente fornece a motivação
e a escolha do cenário. A diferença está em oferecer uma cena
tridimensional navegável, com água animada, em vez de mapas bidimensionais
estáticos, sob a hipótese de que a dinâmica visual da água reforça a
percepção de risco por parte do observador.

\begin{table}[t]
\caption{Comparação qualitativa com trabalhos relacionados}
\label{tab:comparacao}
\centering
\begin{tabular}{@{}p{2.5cm}p{2.0cm}p{1.6cm}p{1.5cm}@{}}
\toprule
\textbf{Trabalho} & \textbf{Técnica de LOD} & \textbf{Água dinâmica} & \textbf{Contexto} \\
\midrule
Clasen e Hege (2006) & Clipmap esférico & Não & Planeta \\
Strugar (2009) & CDLOD (quadtree) & Não & Terreno \\
Welch (2019) & Clipmap & Sim & Oceano \\
DiPaola, Bhardwaj e Sam (2023) & N/A (SIG) & Não & Costa real \\
\textbf{Proposta} & CDLOD adaptado & Sim & Costa (erosão) \\
\bottomrule
\end{tabular}
\fonte{Elaborado pelos autores (2026).}
\end{table}

Conforme a Tabela~\ref{tab:comparacao}, a contribuição deste trabalho não é
uma técnica inédita, mas a integração de técnicas de renderização de
terreno e de simulação de água, tratadas em trabalhos distintos da literatura, a um cenário de análise ambiental, aliada à
determinação experimental do ponto a partir do qual o LOD se torna
necessário (Seção~\ref{sec:resultados}).

\section{Metodologia e Modelagem Computacional}
\label{sec:metod}

A metodologia foi estruturada em um \emph{pipeline} modular de seis
estágios: a) aquisição e geração dos dados de elevação; b) estruturação do
terreno em \emph{quadtree} e nível de detalhe; c) modelagem visual do
terreno; d) modelo de câmera de navegação; e) simulação da superfície
oceânica; e f) ambiente de desenvolvimento. O protocolo da avaliação experimental decorre
diretamente dessa metodologia e é apresentado na
Seção~\ref{sec:resultados}.

\subsection{Aquisição e Geração dos Dados de Elevação}
\label{sub:aquisicao}

A área de estudo compreende a Baixada Santista (municípios de Santos, São
Vicente e Guarujá, SP), delimitada pela caixa geográfica de latitude
$-23{,}70$ a $-24{,}10$ e longitude $-46{,}60$ a $-46{,}05$. O recorte foi
escolhido por reunir, em uma única área, relevo montanhoso (Serra do Mar),
planície costeira baixa e o sistema estuarino de Santos e São Vicente
(região onde rios desembocam no mar).

O relevo é representado por um mapa de elevação (\emph{heightmap}), isto é,
uma matriz $H \in \mathbb{R}^{n \times n}$, com $n = 1024$, na qual cada
célula $H(x,z)$ armazena a altura do terreno em unidades de mundo
(Equação~\eqref{eq:escala}). Consideram-se dois métodos
de obtenção de $H$, sendo o primeiro a fonte primária da implementação de
referência.

\subsubsection{Dados reais (fonte primária)}
O terreno é obtido a partir de \emph{tiles} de elevação do serviço
\emph{AWS Terrain Tiles}, distribuído publicamente sem exigência de chave
de API, no formato \emph{terrarium}. Nesse formato, a altitude real, em
metros, é codificada nos três canais de cor de uma imagem PNG, segundo a
Equação~\eqref{eq:terrarium}:
\begin{equation}
h_{\text{real}} = (R \cdot 256 + G + B/256) - 32768
\label{eq:terrarium}
\end{equation}
em que $R$, $G$ e $B$ são os valores dos canais de cor do \emph{tile}
(0 a 255) e $h_{\text{real}}$ é a elevação decodificada, em metros. Os dados
subjacentes combinam observações orbitais de radar (\emph{Shuttle Radar Topography
Mission}, SRTM, para a porção emersa) e batimetria da \emph{General
Bathymetric Chart of the Oceans} (GEBCO, para o leito oceânico), o que permite
representar relevo terrestre e profundidade marinha em uma única matriz
contínua.

Como os \emph{tiles} brutos apresentam ruído pontual, característico de
dados de sensoriamento remoto, aplica-se um pós-processamento composto por
três etapas: a) filtro de mediana $3 \times 3$, em duas passadas, para
remover valores isolados espúrios (\emph{despiking}) sem suavizar
excessivamente as feições do relevo; b) reamostragem por interpolação
bilinear para a resolução-alvo de $1024 \times 1024$ células; e c)
imposição de um piso mínimo de elevação de $-40$ m, que evita que
artefatos de batimetria produzam profundidades irrealistas na região de
interesse.

A elevação real, em metros, é convertida para unidades de mundo (u.m.) do
motor gráfico por um fator de escala horizontal e um fator de exagero
vertical, aplicados conjuntamente, conforme a Equação~\eqref{eq:escala}:
\begin{equation}
y_{\text{mundo}} = \frac{h_{\text{real}}}{\sigma} \cdot E
\label{eq:escala}
\end{equation}
em que $\sigma = 167$ m/u.m. é a escala do mundo (um terreno de 300 u.m. de
lado corresponde a cerca de 50 km reais, cobrindo a extensão da área de
estudo) e $E = 6$ é o fator de exagero vertical. O exagero vertical é prática
consolidada em geovisualização e é necessário neste caso porque o relevo
real da Serra do Mar, com amplitude aproximada de 1,2 km, representa uma
fração muito pequena da extensão horizontal da área de estudo, de cerca de
50 km. Sem o fator $E$, o relevo apareceria visualmente quase plano.

\subsubsection{Geração procedural (modo alternativo)}
Como alternativa independente de conectividade de rede, a matriz $H$ pode
ser gerada por ruído coerente do tipo \emph{Simplex noise}, somado em
múltiplas oitavas (\emph{fractal Brownian motion}), conforme a
Equação~\eqref{eq:fbm}:
\begin{equation}
H(x,z) = \frac{1}{\sum_{i=0}^{k-1} p^i} \sum_{i=0}^{k-1} p^i \cdot
\text{Simplex}(f_0 2^i x,\, f_0 2^i z)
\label{eq:fbm}
\end{equation}
com $k = 5$ oitavas e persistência $p = 0{,}5$, em que $f_0 = 2$ é a
frequência base da primeira oitava e $\text{Simplex}(\cdot)$ é uma função de
ruído coerente que devolve valores contínuos entre $-1$ e $1$. Cada oitava
$i$ dobra a frequência ($2^i$) e reduz a contribuição por $p^i$: as oitavas
finas acrescentam detalhes pequenos ao relevo grosseiro das primeiras.

Independentemente da origem dos dados, cada célula de $H$ é usada para o
deslocamento vertical dos vértices de uma grade regular (\emph{displacement
mapping}). A matriz é transferida à GPU como uma textura de ponto flutuante
de canal único (R32F), com um valor de altura por \emph{texel}, lido
diretamente pelo \emph{vertex shader} no estágio descrito na
Subseção~\ref{sub:cdlod}.

\subsection{Estruturação do Terreno em Quadtree e Nível de Detalhe (CDLOD)}
\label{sub:cdlod}

O algoritmo CDLOD (STRUGAR, 2009) resolve o seguinte problema: renderizar
toda a extensão do terreno na resolução máxima é inviável em tempo real, mas
reduzir a resolução de maneira uniforme degrada a qualidade visual nas
regiões próximas ao observador, onde o detalhe é mais perceptível. A solução
consiste em variar a densidade de triângulos em função da distância à
câmera.

\subsubsection{Estrutura de dados}
O terreno é particionado recursivamente em uma \emph{quadtree}. O nó raiz
cobre toda a extensão do terreno e cada nó filho cobre um quadrante da área
do nó pai, até a profundidade máxima $L = 5$, valor implementado e
suficiente para representar a região costeira de interesse com resolução
submétrica nas proximidades do observador. Cada nó armazena a mesma
quantidade fixa de vértices, organizados em uma grade regular de
$N \times N$ pontos ($N = 33$, correspondendo a $32 \times 32$ subdivisões,
ou 2048 triângulos por nó). Assim, um nó na profundidade $\ell$ cobre uma
área de $\left(\frac{W}{2^{\ell}}\right)^2$, sendo $W$ a largura total do
terreno. Nós mais profundos representam áreas menores com a mesma densidade
de vértices, o que resulta em maior resolução por unidade de área.

\subsubsection{Critério de seleção de nível}
Para cada nó de lado $\tau$ (tamanho do bloco na profundidade
correspondente), calcula-se a distância euclidiana $d$, no plano horizontal, entre a posição da
câmera $\vec{c}$ e o ponto mais próximo da borda da caixa delimitadora do
nó. O critério de subdivisão implementado é dado pela
Equação~\eqref{eq:subdivisao}:
\begin{equation}
\text{subdividir se } d < \tau \cdot 2{,}2
\label{eq:subdivisao}
\end{equation}
Ou seja, um nó é subdividido em quatro filhos de maior detalhe sempre que a
câmera estiver mais próxima de sua borda do que 2,2 vezes o tamanho do
próprio bloco. Caso contrário, o nó é renderizado em sua resolução atual.
Como o tamanho do bloco $\tau$ é reduzido à metade a cada nível de
profundidade, o critério equivale a um raio de comutação que dobra a cada
nível, semelhante a um esquema de anéis concêntricos de detalhe centrados na
câmera, sem exigir o armazenamento de distâncias de corte por nível. Na
implementação de referência, esse critério resulta tipicamente em 200 a 230
blocos ativos simultaneamente, contra os $4^5 = 1024$ blocos-folha que
existiriam caso a árvore inteira fosse forçada à profundidade máxima.

\par\smallskip\noindent{\footnotesize Algoritmo 1 \textendash\ Seleção dos blocos da \emph{quadtree} a cada quadro}
{\footnotesize
\begin{verbatim}
SELECIONAR(no, c, k, saida)
  d <- distancia_horizontal(c, borda de no)
  se no tem filhos e d < k * no.lado:
      para cada filho de no:
          SELECIONAR(filho, c, k, saida)
  senao:
      adicionar no a saida
\end{verbatim}}
\fonte{Elaborado pelos autores (2026), com base no código-fonte (\texttt{lod.py}).}

\subsubsection{Eliminação de descontinuidades (\emph{cracks})}
Blocos adjacentes renderizados em níveis de detalhe diferentes possuem
densidades de vértices distintas na borda compartilhada, o que pode gerar
fendas visuais (\emph{cracks}) caso os vértices das duas malhas não
coincidam exatamente na fronteira. O CDLOD trata o problema por meio do
\emph{geomorphing}: na faixa de transição de cada bloco, os vértices de
índice ímpar da grade deslizam progressivamente para a posição dos vértices
de índice par correspondentes. A faixa de transição de um bloco começa a
55\% do caminho entre o seu alcance de seleção, $r = 2{,}2\,\tau$
(Equação~\eqref{eq:subdivisao}), e o alcance do bloco-pai, $2r$, e termina
neste último. Nesse intervalo, o fator de transição de cada vértice,
calculado a partir de sua distância à câmera, cresce linearmente de 0 a 1:
com fator 0 o vértice permanece na posição própria e com fator 1 coincide
com a grade mais grossa do bloco-pai. Isso elimina tanto as fendas quanto o salto
abrupto (\emph{popping}) que ocorreria quando a câmera se move e um bloco
muda de nível de detalhe.

\par\smallskip\noindent{\footnotesize Algoritmo 2 \textendash\ Transição do \emph{geomorphing} por vértice (\emph{vertex shader})}
{\footnotesize
\begin{verbatim}
d <- distancia(c, vertice)
f <- limitar((d - d_ini)/(d_fim - d_ini), 0, 1)
g <- g - f * resto(g)   // vertice impar -> par
y <- amostrar(heightmap, g)
\end{verbatim}}
\fonte{Elaborado pelos autores (2026), com base no código-fonte (\texttt{terrain.vert}).}

\subsection{Modelagem Visual do Terreno}
\label{sub:visual}

A cor de cada ponto do terreno é calculada no \emph{fragment shader} a
partir de quatro componentes: a cor-base por altitude e inclinação, a
iluminação difusa, a tonalidade das áreas submersas e a névoa de distância.

A cor-base depende da altitude normalizada $\hat{h} = h / h_{\max}$, em que
$h_{\max}$ é o maior valor absoluto de altura do terreno. Por interpolação linear, ela passa de
areia de restinga ($\hat{h} < 0{,}04$) para mata da planície e da baixa
encosta, depois para mata densa na encosta da serra ($\hat{h} < 0{,}55$) e,
por fim, para rocha exposta nos cumes. Além disso, a inclinação do terreno
mistura rocha à cor-base em até 55\%, em proporção a
$\min(1,\ 2{,}2\,(1 - n_y))$, em que $n_y$ é a componente vertical da normal, de modo que os paredões íngremes da
escarpa apareçam como rocha e não como mata. A normal é calculada no
\emph{vertex shader} por diferenças finitas, a partir de quatro amostras do
\emph{heightmap} vizinhas ao ponto, conforme a Equação~\eqref{eq:normal}:
\begin{equation}
\vec{n} = \frac{(h_E - h_D,\ 2\varepsilon,\ h_B - h_C)}{\left\lVert (h_E - h_D,\ 2\varepsilon,\ h_B - h_C) \right\rVert}
\label{eq:normal}
\end{equation}
em que $h_E$, $h_D$, $h_B$ e $h_C$ são as alturas amostradas à esquerda, à
direita, abaixo e acima do ponto, a uma distância $\varepsilon$ de 2,5 texels,
valor que suaviza o ruído do \emph{heightmap}.

A iluminação difusa segue a lei de Lambert, com uma parcela de luz ambiente
para que nenhuma região fique totalmente escura, conforme a
Equação~\eqref{eq:lambert}:
\begin{equation}
\text{cor} = \text{cor}_{\text{base}} \cdot \big(0{,}25 + 0{,}85\,\max(\vec{n} \cdot \vec{l},\, 0)\big)
\label{eq:lambert}
\end{equation}
em que $\vec{l}$ é o vetor unitário que aponta do ponto para o Sol.

Os pontos com altura inferior ao nível do mar $y_0$ têm a cor misturada, em
65\%, a um tom azul-esverdeado que escurece com a profundidade. É esse
efeito que torna visível a área inundada nas capturas da
Figura~\ref{fig:capturas}. Por fim, a partir de 120 u.m. da câmera, a cor é
misturada à cor do céu, em proporção crescente até o máximo de 75\%,
atingido a 315 u.m., o que
dá noção de escala e esconde o limite distante do terreno.

\begin{figure*}[t]
\centering
\caption{Capturas da aplicação, a partir da mesma posição de câmera: (a)
nível do mar em 0 m, com as feições identificadas; (b) nível do mar em +10 m,
cenário extremo usado apenas para demonstração; (c) blocos da
\emph{quadtree} coloridos por profundidade, do azul (blocos grandes e
distantes) ao vermelho (blocos pequenos e próximos)}
\label{fig:capturas}
\includegraphics[width=\textwidth]{capturas_sistema.png}
\fonte{Elaborado pelos autores (2026).}
\end{figure*}

\subsection{Modelo de Câmera de Navegação Livre}
\label{sub:camera}

A navegação pela cena é realizada por uma câmera de movimento livre
(\emph{free-fly camera}), cuja orientação é parametrizada pelos ângulos de
Euler \emph{yaw} ($\psi$, rotação em torno do eixo vertical) e \emph{pitch}
($\theta$, rotação em torno do eixo horizontal lateral). O vetor que indica
para onde a câmera aponta (\emph{front}) é derivado desses ângulos pela
Equação~\eqref{eq:front}:
\begin{equation}
\vec{f} = \big(\cos\psi\cos\theta,\ \sin\theta,\ \sin\psi\cos\theta\big)
\label{eq:front}
\end{equation}

A partir de $\vec{f}$ e do vetor global ``para cima''
$\vec{u}_{\text{mundo}} = (0,1,0)$, calculam-se os vetores ortogonais
\emph{right} ($\vec{r} = \vec{f} \times \vec{u}_{\text{mundo}}$,
normalizado) e \emph{up} ($\vec{u} = \vec{r} \times \vec{f}$), que formam a
base local da câmera. Esses três vetores, associados à posição $\vec{c}$,
constroem a matriz de visualização (\emph{view matrix}) por meio da
transformação \emph{look-at}, que expressa a cena no referencial da câmera.

A cada quadro, as entradas de teclado incrementam $\vec{c}$ na direção de
$\vec{f}$, $\vec{r}$ ou $\vec{u}$, com velocidade de 40 unidades de mundo
por segundo no modo normal e de 120 no modo acelerado (\emph{turbo}), ambas
escaladas pelo tempo decorrido entre quadros (\emph{delta time}) para
garantir independência da taxa de quadros. O deslocamento do mouse
incrementa $\psi$ e $\theta$, com $\theta$ restrito ao intervalo
$(-89^{\circ}, 89^{\circ})$ para evitar a inversão do eixo vertical da
câmera (\emph{gimbal flip}). A matriz de projeção perspectiva usa campo de
visão vertical de $60^{\circ}$ e planos de recorte próximo e distante em 0,1
e 2000 unidades de mundo, respectivamente, intervalo dimensionado para conter
toda a extensão do terreno descrita na Subseção~\ref{sub:aquisicao} sem
introduzir imprecisão numérica no \emph{depth buffer}.

Além da navegação, $\vec{c}$ é a entrada direta do critério de distância da
Subseção~\ref{sub:cdlod}: a cada quadro, a mesma posição de câmera usada
para construir a matriz de visualização é reutilizada no cálculo de $d$ para
cada nó da \emph{quadtree}, o que evita manter duas representações
distintas da posição do observador.

\subsection{Simulação da Superfície Oceânica}
\label{sub:agua}

A superfície da água é modelada pelo modelo de ondas de Gerstner, também
chamado trocoidal, na formulação de Fournier e Reeves (1986). Diferentemente
de uma onda senoidal simples, que desloca os pontos da superfície apenas na
vertical, o modelo de Gerstner desloca cada ponto também na direção
horizontal de propagação da onda. As partículas da superfície descrevem,
assim, trajetórias elípticas, o que produz o perfil característico de
cristas estreitas e pontiagudas e vales largos e suaves das ondas oceânicas
reais. Para um ponto de referência $(x_0, z_0)$ no plano da água em repouso,
a posição deslocada por uma única componente de onda é dada pelas
Equações~\eqref{eq:gerstner-x} a~\eqref{eq:gerstner-y}:
\begin{align}
x &= x_0 + Q A \hat{D}_x \cos(\vec{k} \cdot \vec{x}_0 - \omega t) \label{eq:gerstner-x}\\
z &= z_0 + Q A \hat{D}_z \cos(\vec{k} \cdot \vec{x}_0 - \omega t) \label{eq:gerstner-z}\\
y &= A \sin(\vec{k} \cdot \vec{x}_0 - \omega t) \label{eq:gerstner-y}
\end{align}
em que $A$ é a amplitude da onda, $\vec{k}$ é o vetor de onda (de módulo
$2\pi/\lambda$, sendo $\lambda$ o comprimento de onda, e direção de
propagação $\hat{D}$), $\omega$ é a frequência angular, $t$ é o tempo e
$Q \in [0,1]$ controla a nitidez das cristas (\emph{steepness}). Valores
próximos de 1 produzem cristas mais pontiagudas, ao custo de possíveis
autointerseções da malha caso o limite seja excedido.

A superfície final é obtida pela soma de $m = 6$ componentes de onda, que
variam de ondulação longa a ondulação curta, com amplitudes
$A_i \in \{1{,}10;\, 0{,}90;\, 0{,}50;\, 0{,}28;\, 0{,}14;\, 0{,}07\}$ e
comprimentos de onda $\lambda_i \in \{70;\, 48;\, 26;\, 14;\, 7;\, 4\}$
unidades de mundo. A nitidez $Q_i$ decresce de 0,60 a 0,28 à medida que a
frequência aumenta, restrição necessária para evitar a autointerseção da
malha nas componentes de alta frequência, cujo menor comprimento de onda
implica menor tolerância geométrica a cristas pronunciadas. A superposição
evita o padrão artificial e repetitivo de uma única componente e aproxima o
aspecto irregular de uma superfície oceânica real, ao custo de $m$
avaliações da Equação~\eqref{eq:gerstner-y} por vértice a cada quadro.

A normal de cada ponto da superfície, necessária para a iluminação, é obtida
analiticamente, a partir das derivadas parciais da soma das $m$ componentes
em relação a $x_0$ e $z_0$, e não por diferenças finitas, o que evita o
custo de amostrar vértices vizinhos. O modelo de iluminação combina um termo
difuso lambertiano, um termo especular de Blinn-Phong, um termo de
reflectância dependente do ângulo de visão, dado pela aproximação de
Fresnel-Schlick, e transparência com fator $\alpha = 0{,}80$, que permite
perceber parcialmente o relevo submerso através da água.

A altura de base $y_0$ da superfície, que é o plano de repouso em torno do
qual as componentes de Gerstner oscilam, é um parâmetro ajustável do
sistema. Ela pode ser controlada de forma discreta, por quatro cenários
predefinidos (0, +0,5, +1 e +2 metros sobre o nível de referência atual), ou
de forma contínua, a uma taxa de 0,8 m/s enquanto o comando correspondente
é mantido pressionado. A submersão de um ponto do terreno é determinada
pela comparação direta entre a elevação do terreno nessa posição
(Equação~\eqref{eq:escala}) e $y_0$: os pontos de elevação inferior a $y_0$
ficam abaixo da superfície da água e, portanto, visualmente submersos, sem
qualquer alteração na geometria do terreno. O valor de $y_0$ é informado em
metros e convertido para unidades de mundo pelo mesmo fator $E/\sigma$ da
Equação~\eqref{eq:escala}; assim, o exagero vertical atua de modo idêntico
sobre o terreno e sobre a água, e a área submersa corresponde à elevação
real, em metros. A Figura~\ref{fig:detalhes}(a) mostra a superfície
resultante: as cristas estreitas das ondas de Gerstner, o brilho especular
do Sol e a transparência, que deixa ver o terreno submerso.

\begin{figure*}[t]
\centering
\caption{Detalhes do sistema: (a) superfície da água vista de perto, com
as cristas das ondas de Gerstner, o brilho especular do Sol e a
transparência sobre o terreno submerso; (b) interface, com a legenda de
comandos, o painel de status (nível do mar em +2,00 m, modo de detalhe e
tempo de quadro), a bússola e o minimapa}
\label{fig:detalhes}
\includegraphics[width=0.85\textwidth]{detalhes_agua_interface.png}
\fonte{Elaborado pelos autores (2026).}
\end{figure*}

\subsection{Ambiente de Desenvolvimento}
\label{sub:ambiente}

A aplicação foi desenvolvida em Python 3.11. A biblioteca GLFW cria a janela
e captura teclado e mouse. A ModernGL dá acesso ao OpenGL 3.3 (\emph{core
profile}, GLSL versão 330) para compilar \emph{shaders}, alocar \emph{buffers}
na GPU e emitir os comandos de desenho. A PyGLM fornece a álgebra linear das
Equações~\eqref{eq:front} a~\eqref{eq:gerstner-y}. A Pillow decodifica os
\emph{tiles} de elevação (Equação~\eqref{eq:terrarium}) e renderiza o texto da
interface, e a matplotlib gera os gráficos de desempenho. Por serem imagens
PNG convencionais, os \emph{tiles} dispensam bibliotecas geoespaciais, como
GDAL ou \texttt{rasterio}: a decodificação usa apenas Pillow e NumPy.

A seleção dos blocos da \emph{quadtree} (Subseção~\ref{sub:cdlod}) é feita
na unidade central de processamento (CPU) a cada quadro. O deslocamento dos
vértices pelo \emph{heightmap}, a transição do \emph{geomorphing} e as ondas
de Gerstner (Subseção~\ref{sub:agua}) são calculados no \emph{vertex shader},
para que o paralelismo da GPU absorva o custo por vértice. O código é
organizado em oito módulos (aquisição do DEM, geração procedural,
\emph{quadtree}, geometria da água, câmera, cena, interface e laço
principal), além do módulo de teste de desempenho
(Seção~\ref{sec:resultados}). A interface inclui tela introdutória, legenda
de comandos em dois níveis, painel de status (nível do mar, modo de
renderização e tempo de quadro), bússola e minimapa (Figura~\ref{fig:detalhes}b), elementos que buscam
tornar a ferramenta acessível a um público não especializado. A
Tabela~\ref{tab:parametros} sintetiza os parâmetros da implementação de
referência.

\begin{table}[t]
\caption{Parâmetros de configuração da implementação de referência}
\label{tab:parametros}
\centering
\begin{tabular}{@{}p{3.1cm}p{0.9cm}p{2.4cm}@{}}
\toprule
\textbf{Parâmetro} & \textbf{Símbolo} & \textbf{Valor} \\
\midrule
Resolução do \emph{heightmap} & $n$ & $1024 \times 1024$ \\
Exagero vertical & $E$ & $6\times$ \\
Escala do mundo & $\sigma$ & 167 m/u.m. \\
Profundidade máx. da quadtree & $L$ & 5 níveis \\
Resolução de vértices por nó & $N$ & $33 \times 33$ \\
Fator do critério de subdivisão & - & $2{,}2 \times \tau$ \\
Componentes de onda de Gerstner & $m$ & 6 \\
Nitidez das ondas (faixa) & $Q$ & 0,60 a 0,28 \\
Transparência da água & $\alpha$ & 0,80 \\
Campo de visão da câmera & FOV & $60^{\circ}$ \\
Planos de recorte (próximo/distante) & - & 0,1 / 2000 u.m. \\
Velocidade normal / turbo & - & 40 / 120 u.m./s \\
\bottomrule
\end{tabular}
\fonte{Elaborado pelos autores (2026), com base no código-fonte da implementação.}
\end{table}

\section{Resultados e Avaliação de Desempenho}
\label{sec:resultados}

A Figura~\ref{fig:capturas} mostra a aplicação em funcionamento. No painel
(a), com o nível do mar em 0 m, a planície costeira aparece emersa, o
estuário e o oceano aparecem cobertos pela água e a Serra do Mar forma a
escarpa ao fundo. No painel (b), o nível
foi elevado a +10 m, valor escolhido apenas para tornar a inundação
visível na imagem; os cenários de +0,5 a +2 m usam o mesmo mecanismo, com
extensão alagada menor. O painel (c) torna visível o refinamento do CDLOD:
os blocos próximos à câmera têm maior profundidade na \emph{quadtree}. Uma
inspeção visual indica que a forma do terreno, com os canais do sistema
estuarino de Santos e São Vicente, as ilhas próximas e a orientação da
escarpa da Serra do Mar, corresponde à geografia real da região quando
comparada a um mapa de referência, o que sugere que a aquisição e a
decodificação descritas na Subseção~\ref{sub:aquisicao} estão corretas;
não foi feita validação quantitativa.

A Tabela~\ref{tab:area} quantifica a área emersa que fica abaixo do nível
do mar em cada cenário, calculada diretamente sobre o DEM de $1024 \times
1024$ células, em que cada célula cobre cerca de 49 m $\times$ 49 m
(0,0024 km$^2$). Considera-se emersa a célula de elevação maior ou igual a
0 m, o que totaliza 1927 km$^2$. Há 104,6 km$^2$ de células com elevação
exatamente igual a 0 m, que são alagadas assim que o nível passa de zero.
É provável que correspondam a corpos d'água e a áreas costeiras planas do
produto de origem, mas essa hipótese não foi verificada. Descontando essas
células, a área alagada cresce de 14,8 km$^2$ em +0,5 m para 38,9 km$^2$ em
+2 m. Isso explica por que os cenários de +0,5 a +2 m produzem pouca
diferença visual em uma vista oblíqua, como a da Figura~\ref{fig:capturas},
e por que o painel (b) usa +10 m.

\begin{table}[t]
\caption{Área emersa alagada em cada cenário de nível do mar}
\label{tab:area}
\centering
\begin{tabular}{@{}p{1.3cm}p{1.9cm}p{2.0cm}p{2.3cm}@{}}
\toprule
\textbf{Nível} & \textbf{Área alagada (km$^2$)} & \textbf{Parcela da área emersa} & \textbf{Sem as células a 0 m (km$^2$)} \\
\midrule
+0,5 m & 119,4 & 6,2\% & 14,8 \\
+1 m   & 127,7 & 6,6\% & 23,1 \\
+2 m   & 143,5 & 7,4\% & 38,9 \\
+10 m  & 505,0 & 26,2\% & 400,4 \\
\bottomrule
\end{tabular}
\fonte{Elaborado pelos autores (2026), a partir do DEM de 1024 $\times$ 1024
células. Área emersa: células com elevação maior ou igual a 0 m (1927
km$^2$).}
\end{table}

A avaliação de desempenho foi conduzida por um programa de teste
(\emph{benchmark}) dedicado (\texttt{bench.py}), executado sem abrir janela
(modo \emph{headless}). O programa reutiliza a mesma lógica de montagem de
cena (\texttt{scene.py}) da aplicação interativa, de modo que a medição
reflita o custo real de renderização e não uma via de código simplificada.
O protocolo sincronizou a GPU ao final de cada quadro, para que o
\emph{buffering} assíncrono não mascarasse o tempo real de processamento,
descartou os 30 quadros iniciais, de aquecimento (\emph{warm-up}), e mediu
os 150 quadros seguintes.

O experimento variou três fatores: a) a técnica de renderização, com malha
única na resolução total (sem LOD), CDLOD completo e CDLOD sem
\emph{geomorphing} (Subseção~\ref{sub:cdlod}); b) a resolução do
\emph{heightmap}, em cinco níveis (256, 512, 1024, 1536 e 2048); e c) a
posição da câmera, em quatro configurações distintas dentro da cena. Também
foram registrados o consumo de memória do processo (RAM) e o uso de memória
de vídeo (VRAM), por meio da biblioteca \texttt{psutil} e da ferramenta
\texttt{nvidia-smi}, respectivamente. Os testes foram executados em uma
GPU NVIDIA RTX 3060.

A Tabela~\ref{tab:tempos} apresenta o tempo médio de quadro e o percentil
95 (p95), isto é, o tempo abaixo do qual ficam 95\% dos quadros, que indica
a regularidade da animação. Os valores são médias das quatro posições de
câmera. A Figura~\ref{fig:custo} mostra os mesmos dados em função do número
de triângulos da malha completa.

\begin{table}[t]
\caption{Tempo de quadro por resolução do \emph{heightmap}}
\label{tab:tempos}
\centering
\begin{tabular}{@{}p{1.5cm}p{1.7cm}p{2.0cm}p{2.0cm}@{}}
\toprule
\textbf{Resolução} & \textbf{Triângulos sem LOD (milhões)} & \textbf{Sem LOD, média / p95 (ms)} & \textbf{CDLOD, média / p95 (ms)} \\
\midrule
256  & 0,26 & 0,17 / 0,24 & 1,06 / 1,29 \\
512  & 0,65 & 0,40 / 0,67 & 1,21 / 1,70 \\
1024 & 2,23 & 0,74 / 1,27 & 1,06 / 1,35 \\
1536 & 4,85 & 1,18 / 1,70 & 1,07 / 1,33 \\
2048 & 8,52 & 2,05 / 2,75 & 1,10 / 1,42 \\
\bottomrule
\end{tabular}
\fonte{Elaborado pelos autores (2026), com dados do teste de desempenho. Os
triângulos incluem a água. O CDLOD desenha cerca de 0,43 milhão de
triângulos por quadro em todas as resoluções.}
\end{table}

\begin{figure}[t]
\centering
\caption{Tempo de quadro em função do número de triângulos da malha
completa, para a média e o p95. As linhas pontilhadas marcam os
cruzamentos estimados por interpolação linear}
\label{fig:custo}
\includegraphics[width=\columnwidth]{custo_vs_carga.png}
\fonte{Elaborado pelos autores (2026), com dados do teste de desempenho.}
\end{figure}

Com o CDLOD, o tempo médio de quadro permanece entre 1,06 e 1,21 ms em todas
as resoluções, o que corresponde a uma carga estável de cerca de 0,43 milhão
de triângulos por quadro (200 a 230 blocos, Subseção~\ref{sub:cdlod}). Sem
LOD, o tempo cresce com o tamanho da malha, que chega a 8,5 milhões de
triângulos na resolução 2048, contra 2,2 milhões na resolução 1024.

Os valores medidos delimitam o ponto de virada entre as resoluções 1024 e
1536, isto é, entre 2,2 e 4,9 milhões de triângulos. Por interpolação linear
entre esses dois pontos, as curvas de tempo médio se cruzam em cerca de 4,2
milhões de triângulos e as de p95, em cerca de 2,7 milhões
(Figura~\ref{fig:custo}). Abaixo desse ponto, a malha sem LOD é mais rápida,
pois o CDLOD tem uma sobrecarga fixa, a de percorrer a \emph{quadtree} e
emitir cerca de 200 chamadas de desenho a cada quadro, que não se paga em
cargas pequenas. Como o p95 representa o pior caso de regularidade, a
interatividade se perde na malha completa antes do que o tempo médio sugere.
O resultado valida a hipótese da Seção~\ref{sec:intro}: o CDLOD não é
vantajoso de forma incondicional, mas passa a ser indispensável a partir de
uma carga geométrica identificável. Em uma GPU de menor capacidade, espera-se que o ponto de virada ocorra em resolução
menor, o que não foi medido (Seção~\ref{sec:conclusao}). O custo do
\emph{geomorphing}, dado pela diferença entre as condições ``CDLOD'' e
``CDLOD sem \emph{morphing}'', foi de no máximo 0,06 ms, dentro da margem de
ruído; logo, a eliminação de \emph{cracks} e de \emph{popping} tem custo de
desempenho desprezível.

Quanto ao consumo de recursos, a memória do processo (RAM) cresceu de cerca
de 109 MB, na resolução 256, para 256 MB, na resolução 2048. A memória de
vídeo (VRAM) cresceu de cerca de 620 MB para 805 MB no mesmo intervalo; esse
valor é o uso total da GPU informado pelo \texttt{nvidia-smi}, que inclui
outros processos do sistema. O crescimento acompanha a resolução do dado de
entrada (o tamanho da textura de altura e dos \emph{buffers}
correspondentes) e não o modo de renderização: nas condições avaliadas, a
malha completa e as duas variantes do CDLOD compartilham, em uma mesma
execução, os mesmos recursos residentes na GPU. Portanto, o experimento não
isola o ganho de memória atribuível ao LOD, o que fica como trabalho futuro
(Seção~\ref{sec:conclusao}).

\section{Conclusão e Trabalhos Futuros}
\label{sec:conclusao}

Este artigo apresentou um sistema interativo que exibe em três dimensões o
relevo real da Baixada Santista e uma superfície de água animada, cuja
altura pode ser elevada em 0, $+0{,}5$, $+1$ e $+2$ metros, permitindo
observar quais áreas ficariam submersas em cada cenário de subida do mar.

O principal resultado vem do teste de desempenho. Com o CDLOD, o tempo para
desenhar um quadro da cena completa (terreno e água) ficou próximo de 1 ms
em todas as resoluções de terreno testadas. Sem LOD, esse tempo cresce
junto com o tamanho da malha. As duas curvas se cruzam entre 2,7 e 4,2
milhões de triângulos, conforme se use o percentil 95 ou a média: abaixo
desse valor, desenhar a malha completa é mais rápido; acima dele, o CDLOD
passa a ser indispensável para manter a interatividade. O custo do \emph{geomorphing} mostrou-se
desprezível. A eficácia do CDLOD já é conhecida na literatura (STRUGAR,
2009); o que este trabalho acrescenta é a medição desse ponto de virada em
um cenário aplicado de inundação costeira. Assim, o objetivo proposto foi
atingido e a hipótese inicial, de que o nível de detalhe adaptativo é o
mecanismo que viabiliza a visualização tridimensional de inundação costeira
em tempo real, foi confirmada.

Algumas limitações devem ser consideradas. Primeiro, a implementação do
CDLOD é simplificada: não descarta os blocos fora do campo de visão da
câmera (\emph{frustum culling}) e as faixas de transição do
\emph{geomorphing} não seguem exatamente as fórmulas do trabalho original
(STRUGAR, 2009). Segundo, o teste foi feito apenas em uma GPU de computador
de mesa (RTX 3060); em outras GPUs o ponto de virada pode ser diferente. Terceiro, o relevo foi exagerado verticalmente em $6\times$
(Equação~\eqref{eq:escala}) para que a Serra do Mar fosse perceptível na
escala da área, o que é uma distorção deliberada, comum em geovisualização,
mas que precisa ser declarada. Quarto, a água não interage com o terreno
(não há espuma na costa nem refração da luz), e a medição de memória não
separa o efeito do LOD, pois os três modos testados compartilham os mesmos
recursos na GPU. Quinto, a exatidão vertical dos dados de origem não foi
validada: se ela for da ordem dos próprios cenários (+0,5 a +2 m), as áreas
alagadas devem ser lidas como demonstração do método, e não como previsão.
Sexto, a análise estatística se limita à média e ao percentil 95 de uma única
série de medições, sem intervalos de confiança, e o ponto de virada foi
estimado por interpolação linear entre duas resoluções; assim, os valores de
2,7 e 4,2 milhões de triângulos são estimativas, e não medidas diretas.

Como trabalhos futuros, propõem-se seis passos, cada um respondendo a uma
limitação apontada: a) repetir o teste de desempenho em outras GPUs, inclusive integradas, para confirmar o ponto de virada nelas e
medir a memória de cada modo separadamente; b) implementar o
\emph{frustum culling}, de modo a não desenhar os blocos atrás da câmera e
reduzir ainda mais o custo por quadro; c) acrescentar espuma e arrebentação
na linha de costa, para tornar a água mais realista; d) apresentar a
ferramenta a especialistas em gestão costeira, a fim de avaliar se ela de
fato auxilia a comunicar o risco de inundação; e) validar a exatidão
vertical do DEM contra uma fonte de maior precisão, avaliar a sensibilidade
das áreas alagadas ao erro vertical e tratar as células de elevação nula; e
f) repetir o teste de desempenho em execuções independentes, com mais
resoluções entre 1024 e 1536, para obter intervalos de confiança e medir
diretamente o ponto de virada.

\section*{Referências}
\begingroup
\footnotesize
\raggedright
\setlength{\parindent}{0pt}
\setlength{\parskip}{4pt}

BROWN, I.; JUDE, S.; KOUKOULAS, S.; NICHOLLS, R.; DICKSON, M.; WALKDEN,
M. Dynamic simulation and visualisation of coastal erosion. \textbf{Computers,
Environment and Urban Systems}, v. 30, n. 6, p. 840-860, 2006.

CLASEN, M.; HEGE, H.-C. Terrain rendering using spherical clipmaps. In:
EUROVIS 2006. \textbf{Proceedings} [...]. [S. l.]: [s. n.], 2006. p. 91-98.

DIPAOLA, F.; BHARDWAJ, A.; SAM, L. Generating an interactive online map of
future sea level rise along the North Shore of Vancouver: methods and
insights on enabling geovisualisation for coastal communities. \textbf{arXiv},
[S. l.], arXiv:2304.07469, 2023. Disponível em:
\url{https://arxiv.org/abs/2304.07469}. Acesso em: 26 set. 2026.

FOURNIER, A.; REEVES, W. T. A simple model of ocean waves. \textbf{ACM
SIGGRAPH Computer Graphics}, v. 20, n. 4, p. 75-84, 1986.

STRUGAR, F. Continuous distance-dependent level of detail for rendering
heightmaps. \textbf{Journal of Graphics, GPU, and Game Tools}, v. 14, n. 4,
p. 57-74, 2009.

TANNER, C. C.; MIGDAL, C. J.; JONES, M. T. The clipmap: a virtual mipmap.
In: ANNUAL CONFERENCE ON COMPUTER GRAPHICS AND INTERACTIVE TECHNIQUES, 25.,
1998, [S. l.]. \textbf{Proceedings} [...]. [S. l.]: ACM, 1998. p. 151-158.

WELCH, R. \textbf{Real-time ocean simulation}. 2019. Trabalho de Conclusão
de Curso (Bacharelado em Ciência da Computação) -- [S. l.], 2019.
Disponível em: \url{https://ryanwelch.co.uk/project/ocean-simulation/}.
Acesso em: 26 set. 2026.
\endgroup

\end{document}
